\chapter{Bilan du Projet et Retour d'Expérience}
\label{chap:bilan}

\section{Introduction}
Ce quatrième chapitre dresse le bilan d'ingénierie, méthodologique et professionnel de ce stage de fin d'études mené au sein de la société suisse \textbf{\companyName} et de sa filiale \textbf{\companySubsidiary} (Groupe \companyGroup). Il permet de mesurer l'impact des réalisations techniques, d'analyser les compétences développées, de documenter les défis majeurs surmontés et d'évaluer la valeur ajoutée apportée au système d'information de l'entreprise et de ses clients industriels.

\section{Apports et Compétences Développées}

\subsection{Apports Techniques et d'Ingénierie Logicielle}
La conception, l'implémentation et la validation du module \textbf{\odooModule} ont permis de consolider des compétences avancées dans plusieurs domaines de pointe :

\begin{itemize}[leftmargin=1.5cm, label=\textcolor{odooPurple}{\small\textbullet}]
    \item \textbf{Maîtrise Approfondie de l'ORM Odoo 18 \& Métaprogrammation Python} : Compréhension intime du cycle de vie des transactions de l'ORM, manipulation avancée des contextes, héritage de modèles (\textit{Classical \& Delegation Inheritance}) et écriture de requêtes groupées performantes via l'API \texttt{\_read\_group} optimisée.
    \item \textbf{Ingénierie de l'IA Générative et Couche BYO-LLM} : Conception d'une couche d'abstraction multi-fournisseurs (\texttt{LLMProvider}) intégrant Google Gemini (\texttt{gemini-2.5-flash}) et Cerebras. Maîtrise des techniques de \textit{prompt engineering}, d'extraction de données à schéma JSON strict, d'isolation stricte des outils et d'implémentation de flux de confirmation explicite avant écriture.
    \item \textbf{Validation E2E de Navigation et Automatisation Navigateur (Playwright)} : Mise en œuvre d'un banc de test automatisé sous Playwright (Chromium) validant les interactions utilisateurs réelles : dépliage des listes groupées Odoo 18, navigation par fil d'Ariane, maintien de session dans les assistants modaux et drilldown sur tableaux de bord.
    \item \textbf{Approche Qualité Multi-Niveaux \& DevOps} : Élaboration d'une stratégie de vérification combinant 138 tests backend Odoo, 46 tests mock LLM, tests en direct des APIs et 32 tests de navigation E2E, le tout orchestré sous conteneurs Docker et environnement WSL.
\end{itemize}

\subsection{Apports Méthodologiques et Organisationnels}
Sur le plan de la gestion de projet et de la rigueur d'ingénierie logicielle :
\begin{itemize}[leftmargin=1.5cm, label=\textcolor{odooPurple}{\small\textbullet}]
    \item \textbf{Conduite de Projet Agile/Scrum} : Découpage itératif en quatre sprints de deux semaines avec livrables testés et alignement régulier avec l'encadrant industriel chez IT-Koncept.
    \item \textbf{Gestion Proactive de la Dette Technique} : Identification précoce des 15 ruptures de compatibilité et anomalies d'interface Odoo 18, prévenant les régressions en production.
    \item \textbf{Traçabilité et Preuves Matérielles} : Culture de la preuve systématique par l'automatisation de captures d'écrans, de rapports d'exécution HTML et de vidéos de démonstration de bout en bout.
\end{itemize}

\subsection{Apports Humains et Professionnels}
L'immersion au sein de IT-Koncept SA a enrichi ma posture d'ingénieur logiciel :
\begin{itemize}[leftmargin=1.5cm, label=\textcolor{odooPurple}{\small\textbullet}]
    \item \textbf{Compréhension des Exigences des PME Suisses} : Intégration des contraintes réelles de précision, de haute disponibilité et de conformité légale (nLPD) des clients de Suisse romande.
    \item \textbf{Autonomie et Capacité de Résolution de Problèmes Complexes} : Capacité à diagnostiquer des défaillances subtiles (filtrage de lecture par les règles de sécurité, fermeture inopinée de fenêtres modales, comportements de rendu de listes groupées).
\end{itemize}

\section{Défis Majeurs Rencontrés et Solutions Apportées}
La réalisation de ce projet d'envergure a confronté l'équipe à plusieurs défis techniques stimulants, résumés dans le tableau \ref{tab:challenges_solutions}.

\begin{table}[!ht]
\centering
\small
\renewcommand{\arraystretch}{1.18}
\begin{tabularx}{\textwidth}{p{4.0cm} p{4.5cm} X}
\toprule
\textbf{Défi Technique Rencontré} & \textbf{Cause Sous-Jacente} & \textbf{Solution Technique Implémentée} \\
\midrule
\textbf{Rupture d'API \texttt{\_read\_group} Odoo 18} & L'ORM v18 retourne des tuples avec agrégat \texttt{:count} au lieu de dictionnaires & Développement du helper versionné \texttt{read\_group\_counts()} dans \texttt{odoo\_compat.py} \\
\midrule
\textbf{Erreur 500 sur l'action des ordres} & Présence du mode \texttt{gantt} non supporté en Community (\texttt{web\_gantt} absent) & Suppression du mode \texttt{gantt} et adaptation vers les vues liste, calendrier et planning \\
\midrule
\textbf{Masquage des pièces détachées pour l'admin} & Règle \textit{write-own} omettant \texttt{perm\_read} (Odoo applique \texttt{True} par défaut) & Définition explicite de \texttt{perm\_read="0"} sur les règles de pièces détachées \\
\midrule
\textbf{Fermeture prématurée des wizards modaux} & Les méthodes retournaient \texttt{True}, interprété par Odoo comme « fermer la fenêtre » & Retour d'une action de rechargement (\textit{reload action}) maintenant la vue ouverte \\
\midrule
\textbf{Détection d'équipements en langage naturel} & La recherche initiale exigeait une correspondance exacte de phrase complète & Implémentation d'une recherche par sous-chaîne sur les noms, codes et séries \\
\midrule
\textbf{Arrêt inopiné du conteneur en environnement WSL} & Mise en veille d'inactivité de l'instance WSL et conflits de processus hôtes & Conception du script \texttt{e2e-run.sh} avec volume nommé et surveillance active \\
\bottomrule
\end{tabularx}
\caption{Synthèse des défis techniques majeurs et des solutions d'ingénierie appliquées}
\label{tab:challenges_solutions}
\end{table}

\section{Conclusion}
Ce bilan témoigne de l'aboutissement complet des objectifs fixés. L'alliance entre rigueur méthodologique, excellence technique sur l'ORM Odoo 18 et automatisation poussée de la validation de navigation sous Playwright a permis de délivrer un produit robuste, sécurisé et directement exploitable au sein du catalogue de solutions de \textbf{\companyName}.
